\documentclass[10pt]{article}

\usepackage[utf8]{inputenc}

\usepackage[T1]{fontenc}

\usepackage{amsmath}

\usepackage{amsfonts}

\usepackage{amssymb}

\usepackage[version=4]{mhchem}

\usepackage{stmaryrd}

\usepackage{graphicx}

\usepackage[export]{adjustbox}

\graphicspath{ {./images/} }



\begin{document}

решения классификационных задач. Наиболее употребительны следующие модели: модели, построенные с использованием принципа разделения, задающие класс поверхностей, разделяющих образы; модели, построенные на принципе потенциалов; модели вычисления оценок (голосования); структурные модели; статистич. модели.



На уровне модели ставится задача отыскания экстремального по качеству алгоритма распознавания (модель вычисления оценок). О качестве работы распознающего алгоритма обычно судят по результатам работы алгоритма на нек-ром тестовом наборе объектов (контрольная последовательность), для к-рого исследователю априори известна достоверная классификация. При построении общей теории распознающих алгоритмов наиболее полные результаты получены в рамках алгебраич. подхода. Распознающий алгоритм представляется в виде произведения распознающего оператора и решающего правила. Введение над распознающими операторами операций сложения, умножения, умножения на скаляр позволяет доказать существование в рамках нек-рого алгебраич. расширения исходного набора распознающих операторов такого распознающего алгоритма, к-рый обладает әкстремальным качеством на любой контрольной последовательности.



$\mathrm{K}$ задачам P. о. относятся также задачи минимизации описания исходных объектов, выделения информативных признаков.



Лum.: [1] Ж у р а в л е в Ю. И., “Проблемы кибернетики", 1978, в. 33, с. $5-68 ;$ [2] А й 3 е р м а н М. А., Б р а в е pм а н Э. М., Р о з о н о р Л. И., Метод потенциальных функций в теории обучения машин, М., 1970; [3] В а п н и к В. Н.,



\includegraphics[max width=\textwidth, center]{2023_07_08_c3a4b8e7db113a73ccb6g-1}

$\begin{array}{ll}\text { разов, пер. с англ., М., } 1977 . & \text { П. П. Кольцов. }\end{array}$



РАСПРЕДЕЛЕНИЕ - то же, что обобщенная функuur.



F-РАСПРЕДЕЛЕНИЕ - см. Фищера F-распределеrue.



$\boldsymbol{t}$-РАСПРЕДЕЛЕНИЕ - см. Стьюдента распределение.



T्2-РАСПРЕДЕЛЕНИЕ - см. Хотеллинга $T^{2}$-pacnpeделение.



W-РАСПРЕДЕЛЕНИЕ - см. Уищарта распределение.



₹-РАСПРЕДЕЛЕНИЕ - см. Фишера z-распределение .



В-РАСПРЕДЕЛЕНИЕ - см. Бета-распределение .



Г-РАСПРЕДЕЛЕНИЕ - см. Гамма-распределение .



$\chi^{2}$-РАСІРЕДЕЛЕНИЕ - см. “Xu-квадpam» pacnpeделение.



$\omega^{2}$-РАСПРЕДЕЛЕНИЕ - см. «Олега-квадрат» pacпределение.



РАСПРЕДЕЛЕНИЕ ВЕРОЯТНОСТЕИ - одНо ИЗ основных понятий вероятностей теории и математической статистики. При современном подходе в качестве математич. модели изучаемого случайного явления берется соответствующее вероятностное пространство $\{\Omega, S, P\}$, где $\Omega$ - множество әлементарных событий, $S$ - выделенная в $\Omega \sigma$-алгебра подмножеств, $P$ - определенная на $S$ мера со свойством $P(\Omega)=1$ (вероятностная мера).



Любую такую меру на $\{\Omega, S\}$ и называют р а с п р ед е л е ни е м в е р я т н о т е й (см. [1]). Однако это определение, являющееся основой аксиоматики А. Н. Колмогорова (1933), при последующем развитии теории оказалось чрезмерно общим и было, с целью исключения нек-рых «патологических» случаев, заменено более ограничительным, напр. требованием "совершенности» вероятностной меры $P$ (см. [2], особенно добавление к англ. переводу [3] этой книги). От Р. в. в функциональных пространствах требуют обычно выполнения нек-рого свойства регулярности, формулируемого как сепарабельность, но допускающее форму- лировку и в других терминах (см Сепарабельный процесс, а также [4]).



Р. в., встречающиеся в большинстве конкретных задач теории вероятностей и математич. статистики, весьма немногочисленны. Все они давно известны и связаны с основными вероятностными схемами [5]. Они описываются или вероятностями отдельных значений (см. Дискретное распределение) или плотностями вероятности (см. Непрерывное распределение). В необходимых случаях для них составлены таблицы [6].



Из этих основных Р. в. одни связаны с повторениями независимых испытаний (см. Биномиальное распределение, Геометрическое распределение, Полиномиальное распределение), а другие - с соответствующими әтой вероятностной схеме предельными закономерностями, возникающими при неограниченном возрастании числа испытаний (см. Норлальное распределение, Пуассона распределение, Арксинуса распределение). Однако эти предельные распределения могут возникать и как точные: в теории случайных процессов (см. Винеровский процесс, Пуассоновский процесс) или как решения уравнений, возникающих в т. н. характеризационных теоремах (см. также Нормальное распределение, Показательное распределение). Равномерное распределение вероятностей, принимаемое обычно как математич. выражение равновозможности соответствующих исходов опыта, также может быть получено как предельное (напр., при приведении по какому-то модулю суммы большого числа случайных величин или каких-либо других случайных величин с достаточно гладкими «расплывающимися» распределениями). Из упомянутых выше основных Р. в. получают другие Р. в. с помощью функциональных преобразований рассматриваемых случайных величин. Так, напр., в математич. статистике из нормально распределенных случайных величин получают величины, имеющие «xи-квадрат» распределение, нецентральное "хи-квадрат» распределение, Стьюдента распределение, Фишера $F$-распределение и др.



Важные классы распределений были открыты в связн с развитием асимптотич. методов теории вероятностей и математич. статистики (см. Предельные теоремы, Устойчивое распределение, Безгранично делимое распределение, «Олега-квадрат» распределение).



C теоретической и прикладной точек зрения важно умение определить понятие близости распределений. Совокупность всех Р. в. на $\{\Omega, S\}$ может быть различными способами превращена в топологич. пространство. При әтом основную роль играет слабая сходимость Р. в. (см. Распределений сходимость). В одномерном и конечномерном случаях основным средством изучения сходимости Р.в. является аппарат характеристических функций.



Часто полное описание Р. в. (напр., при помощи плотности вероятности или распределения бункции) заменяют заданием небольшого числа характеристик. Из этих характеристик наиболее употребительны математическое ожидание (среднее значение), дисперсия, медиана и моменты в одномерном случае. О числовых характеристиках многомерных Р. в. см. Корреляция, Регрессия.



Статистическим аналогом Р. в. является элпирическое распределение. Эмпирич. распределение и его характеристики могут быть использованы для приближенного представления теоретического Р. в. и его характеристик (см. Оченка статистическая). О способах измерения степени согласия эмнирич. распределения с гипотетическим Р. в. см. Статистических гипотез проверка, Непараметрические методы статистики.



Лит.: [1] К о л м о го р о в А. Н., Основные понятия теории вероятностей, 2 изд., М., 1974; [3] Г н е д е н к о Б В., К о л м о г о р о в А. Н., Предельные распределения для сумм

независимых случайных величин, м.- Л., 1949; [3] и х же, Limit distributions for sums of independent random variables,





\end{document}